\section{Halo, background and nuclear-spin inputs}\label{app:E.27}

A limit from a silicon crystal that is valid for every velocity distribution of the halo, not only the standard one, is the solution of a linear program whose columns are the counts that single streams of dark matter would produce in each bin. A crystal kept at sea level for a hundred days before it goes underground carries tritium whose $\beta$ spectrum puts a fraction of an event per kilogram, day and keV below \SI{1}{keV}. And a spin-dependent coupling to unpolarized nuclei produces no coherent phonon at all. These three facts belong to \cref{sec:21.2,sec:21.7,sec:11.6}, and each needs one computation, specified below.

\subsection{The directional halo-independent bound}

\Cref{sec:21.2} writes the expected count in bin $i$ as $\int_{\mathcal{B}}\mathsf{k}_i(\vect{w})\,\dd\Xi(\vect{w})$, linear in the nonnegative measure $\Xi$ of \cref{eq:21.2-stream} on the ball $\mathcal{B} = \{|\vect{w}| \le v_{\mathrm{esc}}\}$ of galactic-frame velocities ($v_{\mathrm{esc}}$, $v_0$ and $\rho_\chi$ as in \cref{sec:1.2}), with the stream kernel $\mathsf{k}_i(\vect{w})$ the count that a single stream of unit $\Xi$-mass at $\vect{w}$ would leave in bin $i$. A crystal fixed in the laboratory sees the stream at velocity $\vect{w} - \vect{v}_E(t)$ and turns with the sidereal rotation. The kernel is then the time average over the bin of the master rate of \cref{sec:8.2} with the crystal's direction-dependent response $S(\vect{q}, \omega; \Lambda)$ and the final-state weights of \cref{sec:8.7}. The largest value of a linear functional $\int\mathsf{f}\,\dd\Xi$ allowed by the measured upper counts $N_i^{\mathrm{up}}$ is the optimum of the semi-infinite program of \cref{eq:C.12-sip}, with nonnegative dual weights $\lambda_i$; it is attained by at most as many streams as there are bins (\cref{sec:21.2}). The functional $\mathsf{f}$ can be a halo integral above a speed in a direction, the rate of a second crystal or of a second orientation of the same crystal, or a modulation amplitude; a signed $\mathsf{f}$ is allowed provided some bin sees every stream on which $\mathsf{f} > 0$. The method extends the stream approach of nuclear-recoil analyses~\cite{feldstein2014new,ibarra2017optimized} to electronic and directional responses.

The computation is a sweep followed by a small linear program. The response must be resolved in the direction of $\vect{q}$, not averaged over angles. The non-self-consistent band calculation behind it (\cref{app:E.1}) therefore runs on the full zone without symmetry reduction, and the library sits at the truncation label $\mathfrak{T}$ of the target's response tables (\cref{app:E.14,app:E.16}). The stream grid has 768 directions on an equal-area pixelization of the sphere~\cite{gorski2005healpix}, 55 galactic-frame speeds in steps of \SI{10}{km/s} up to $v_{\mathrm{esc}}$, and 24 sidereal phases, about $10^6$ kernel columns. For each stream and phase the energy delta function of the master rate fixes $\omega$ on every vector of the $\vect{q}$ grid. Each column is then one pass over about $10^6$ $\vect{q}$ vectors, some $10^{12}$ multiply-adds per mass, 5 to 50 GPU-hours for about 50 masses. The linear program has one constraint per bin and about \num{4e4} variables and is solved by a dual simplex code~\cite{huangfu2018parallelizing} in seconds.

The grid is refined around the few optimal streams by column generation, and the dual vector $\boldsymbol{\lambda}$ must satisfy $\sum_i\lambda_i\mathsf{k}_i \ge \mathsf{f}$ on the refined grid; that inequality certifies the bound, and the result is reported at two grid levels. In the laboratory frame, a kernel that depends only on the speed must reproduce the single-bin result $N_i^{\mathrm{up}}/\inf_{v \ge u_0}[v\,\mathsf{k}_i(v)]$ for the halo integral above $u_0$; the same statement in galactic speeds does not hold, because there the kernel depends on $|\vect{w} - \vect{v}_E(t)|$. And the kernel must be nonnegative for every stream. A response truncated at some order can have negative pieces. A dual vector that satisfies the dual constraint of \cref{eq:C.12-sip} still bounds the objective by weak duality, but the mass bound behind the grid bracket fails for the computed kernel, and the leakage-corrected bound of \cref{sec:21.2-leak} applies instead. The support constraint is imposed in the galactic frame, a per-bin confidence level needs the Bonferroni correction across bins, and slow streams alias on a coarse $\vect{q}$ grid, where the energy delta function falls between grid points.

\subsection{Cosmogenic production}

A nuclide produced at a rate $R_k$ during exposure $k$ of length $t_k$ and decaying with constant $\lambda_{\mathrm{dec}}$ has, after the whole history, the activity $\mathcal{A}_{\mathrm{act}} = \sum_kR_k(1 - \ee^{-\lambda_{\mathrm{dec}}t_k})\,\ee^{-\lambda_{\mathrm{dec}}t'_k}$, where $t'_k$ is the time from the end of exposure $k$ to the measurement; at saturation $\mathcal{A}_{\mathrm{act}} = R$. The sub-keV rate density is $\mathcal{A}_{\mathrm{act}}$ times the fraction of the $\beta$ spectrum below \SI{1}{keV}, per keV, plus the lines of electron capture. L- and M-shell captures produce lines near 1 and \SI{0.1}{keV}~\cite[Sect.~1]{armengaud2017measurement}, so capture nuclides land in the sub-keV region through their L and M fractions. Tritium in silicon is produced at sea level at $124 \pm 25$ atoms per kilogram and day from all cosmic-ray components, $112 \pm 24$ of them from neutrons~\cite{saldanha2020cosmogenic}, has a half-life of 12.3 years, and puts \SI{8.9}{\percent} of its $\beta$ spectrum below \SI{1}{keV}; a hundred days of exposure then give \SI{0.17}{events.kg^{-1}.d^{-1}.keV^{-1}} below \SI{1}{keV}, and saturation gives 11.

What is missing is the sea-level production rate of each relevant nuclide, in atoms per kilogram and day, for every target retained in \cref{sec:21.7} (the survivors) that has no measured or published rate. The semi-empirical cross-sections of ACTIVIA~\cite{back2008activia}, with measured cross-sections where they exist, are folded with the measured sea-level neutron spectrum, scaled to New York City at the midpoint of solar modulation, and with the measured sea-level proton spectra, the inputs of the silicon measurement~\cite{saldanha2020cosmogenic}. For the light nuclides tritium and $^7$Be, where the semi-empirical systematics fail, cross-sections from a Geant4 hadronic list~\cite{agostinelli2003geant4} or from TALYS are folded with the same spectra. The inputs are the composition and isotopic abundances, the spectrum parameterization and the exposure history. The calculation must reproduce the measured rates of silicon~\cite{saldanha2020cosmogenic} and germanium~\cite{armengaud2017measurement} within their errors before any other material is computed. The models disagree most for tritium: in germanium they span 27.7 to 210 nuclei per kilogram and day~\cite[Sect.~6.3]{armengaud2017measurement}, so the light nuclides carry the largest error. Tritium also diffuses out of silicon, which makes its activity in the crystal depend on the thermal history, and the K:L:M capture fractions are needed per nuclide. ACTIVIA takes seconds per material and Geant4 with $10^8$ primaries 1 to 3 CPU-hours, a few tens to a hundred core-hours for the survivor pool.

\subsection{Spin-dependent nucleon inputs}

For unpolarized nuclear spins the coherent one-phonon amplitude of a spin-dependent coupling vanishes, and only the self response $S^{\mathrm{self}}_i$ of the isotopes $i$ that carry a spin $J_{\mathrm{nuc},i}$ survives, \cref{eq:11.6-spinS}~\cite{bloch2025broadband}. For the spin-dependent operator $\mathcal{O}_4$ the rate contracts it with the nuclear coefficient $c_4^{(i)} = [c_4^{(p)}\langle S_p\rangle + c_4^{(n)}\langle S_n\rangle]/J_{\mathrm{nuc},i}$ of \cref{eq:8.1-nuclear}. Here $\langle S_p\rangle$ and $\langle S_n\rangle$ are the expectation values of the proton and neutron spins in the nuclear ground state; the weight is that of the nucleus, not of a free nucleon. The one-phonon term of $S^{\mathrm{self}}_i$ (\cref{sec:11.6-spins}) contains the local density of states $g_i$ of the isotope, \cref{eq:11.6-lifshitz}, which differs from the host density of states of species $\kappa$ near the edges of the optical band. A coherent piece proportional to the squared polarization returns only in a polarizing field.

Two inputs are missing for the survivors. The spin expectation values come from large-scale shell-model calculations: $\langle S_n\rangle = 0.156$ and $\langle S_p\rangle = 0.016$ for $^{29}$Si, and $\langle S_n\rangle = 0.439$ and $\langle S_p\rangle = 0.031$ for $^{73}$Ge, where a second interaction gives 0.450 and 0.006~\cite[Table~III]{klos2013large}. Each odd-$A$ isotope of the survivors needs values from at least two interactions, and their spread is the nuisance. The host density of states $g_{\mathrm{ph},\kappa}$ comes from density-functional perturbation theory (\cref{app:E.4,app:E.17}), interpolated to a fine grid, with the $\vect{Q}$ grid refined until $\langle\omega^{-1}\rangle$ over $g_{\mathrm{ph},\kappa}$ changes by less than \SI{1}{\percent} and an energy bin no wider than a quarter of the threshold; the local $g_i$ of each isotope follows from it by \cref{eq:11.6-lifshitz}. Reusing the coherent one-phonon response with spin weights in place of the scalar couplings gives zero for the wrong reason, because the coherent amplitude averages out. Even--even isotopes have no spin and contribute nothing, so an element's weight is set by its odd-$A$ abundance alone. The density-of-states runs cost 10 to 100 core-hours over the survivors.
